\section{Introduction}
Superconductivity with nonzero center-of-mass pairing momentum provides a route beyond a spatially uniform condensate. Fulde and Ferrell and, independently, Larkin and Ovchinnikov showed that spin- or exchange-split fermionic states can favor finite-momentum pairing and a spatially modulated order parameter~\cite{FuldeFerrell1964,LarkinOvchinnikov1965,CasalbuoniNardulli2004}. The resulting FFLO problem is intrinsically thermodynamic: the preferred state must be selected by comparing competing minima over both the pairing amplitude and momentum, rather than by assigning a phase from a single local observable.

Altermagnets add a distinct setting for this competition. They combine compensated magnetic order with momentum-dependent spin splitting even though the net magnetization vanishes~\cite{SmejkalSinovaJungwirth2022}. Recent model studies have shown that such spin splitting can support zero-field finite-momentum superconductivity and other unconventional pairing states~\cite{ChakrabortyBlackSchaffer2024,HongParkKim2025}. These results motivate an altermagnetic route to finite-$q$ pairing, but they do not determine the phase structure of the one-dimensional model studied here.

Topology adds a second, logically separate requirement. Finite-momentum pairing can coexist with topological superfluidity in spin-orbit-coupled settings~\cite{ZhangYi2013}, while gapped one-dimensional Bogoliubov--de Gennes (BdG) systems in symmetry class D admit a $\mathbb Z_2$ distinction expressible through a Pfaffian invariant~\cite{Kitaev2001,BudichArdonne2013}. The invariant is meaningful only after bulk-gap applicability is established. A topology-resolved FFLO map must therefore first identify the thermodynamic minimum and its numerical reliability, and only then classify trusted gapped FFLO states.

This hierarchy is computationally demanding. In the present problem, every accepted parameter point requires construction and minimization of $F(\Delta,q)$, comparison with the normal reference, adaptive checks of shallow or edge minima, and, where applicable, full-zone spectral and Pfaffian diagnostics. Uniform grids spend many exact calls far from the structures of interest, and their cost grows rapidly when the parameter domain is enlarged. Active-learning methods can instead concentrate experiments or simulations near uncertain phase boundaries and support batched evaluation~\cite{TerayamaEtAl2019,DaiBrussGlotzer2020,LookmanEtAl2019}. For a physical phase diagram, however, an acquisition model should decide where to calculate next, not replace the calculation that defines the phase.

Here we develop a reliability-aware active-learning workflow in which an audited exact BdG oracle remains the sole authority for thermodynamic and topological labels. We first reconstruct the two-dimensional Normal, uniform-superconducting, and FFLO structure while retaining a separate hard-risk numerical layer. On trusted gapped FFLO states, full-zone gap and Pfaffian diagnostics then resolve cFFLO/tFFLO structure with an explicit neighborhood-size sensitivity audit. Finally, we extend the same oracle--scheduler separation to a three-dimensional parameter domain. The extension reveals thermodynamic and topology-aware structure, while its strict topology convergence criteria remain open; it is not presented as a complete three-dimensional topology atlas.
